\chapter{15回授業計画}
\label{app:syllabus}

本書を教科書として半期15回の授業で使用する場合の授業計画を示す。各回90分の授業を想定し、講義と演習のバランスを考慮している。

\section{授業計画表}

\footnotesize
\begin{longtable}{ccp{1.8cm}p{3.3cm}p{1.8cm}}
\toprule
回 & 対応章 & テーマ & 授業の進め方 & 課題 \\
\midrule
\endfirsthead
\toprule
回 & 対応章 & テーマ & 授業の進め方 & 課題 \\
\midrule
\endhead
\midrule
\multicolumn{5}{r}{（次ページに続く）}\\
\endfoot
\bottomrule
\endlastfoot

1 & 第1章 & 研究活動と生成AI入門 & 【講義60分】AI活用の意義と倫理、批判的AI活用の考え方を解説。【演習30分】各自が普段使っているAIツールとその使い方を共有するワークショップ。 & 自分の研究活動におけるAI活用の現状と課題を400字でまとめる \\
\hline
2 & 第2章前半 & プロンプトエンジニアリング基礎 & 【講義40分】ROCF（Role, Output, Context, Format）の解説と実例紹介。【演習50分】同じタスクに対して異なるプロンプトを試し、出力の違いを比較する実践演習。 & 自分の研究テーマに関するプロンプトをROCFで設計し、3パターン作成 \\
\hline
3 & 第2章後半 & プロンプトエンジニアリング実践 & 【講義30分】Few-shot, Chain-of-Thought, 段階的プロンプトの解説。【演習60分】研究テーマに関する実践的なプロンプト作成と改善のサイクルを体験。 & 最も効果的だったプロンプトとその改善過程をレポートにまとめる \\
\hline
4 & 第3章 & AIを活用した文献調査 & 【講義30分】AI時代の文献調査戦略、ハルシネーション対策。【演習60分】自分の研究テーマについて、AIを使って文献サーベイの初期段階を実施。 & 文献調査レポート（先行研究5本以上の要約と比較表） \\
\hline
5 & 第4章 & 英語論文の読解と執筆支援 & 【講義40分】学術英語の構造、AIを使った英文作成・校正の方法。【演習50分】自分の研究のAbstractを日本語→英語で作成する実践。 & 自分の研究の英語Abstract（150--250語）を作成 \\
\hline
6 & 第5章 & プレゼンテーション準備 & 【講義30分】効果的な学会発表の構成、AIを使ったスライド作成。【演習60分】5分間の研究紹介スライド（5--7枚）の構成をAIと共に設計。 & 研究紹介スライドの完成版を提出 \\
\hline
7 & 第6章前半 & プログラミングとAI（基礎） & 【講義30分】AI支援コーディングの原則、GitHub Copilot等の使い方。【演習60分】自分の研究に関連するデータ処理コードをAIと共に作成。 & 作成したコードとAIとのやり取りの記録を提出 \\
\hline
8 & 第6章後半 & プログラミングとAI（応用） & 【講義20分】コードレビュー、テスト、デバッグへのAI活用。【演習40分】前回のコードのリファクタリングとテスト作成。【発表30分】中間発表（3分×12名程度）。 & コードを改善し、README付きで提出 \\
\hline
9 & 第7章 & データ分析・可視化とAI & 【講義30分】AIを使ったデータ分析の計画立案、グラフ作成の原則。【演習60分】自分の研究データ（またはサンプルデータ）を使った分析と可視化の実践。 & データ分析結果のグラフ3枚と各グラフの解釈を記述したレポート \\
\hline
10 & 第8章 & GitHub Webページの作成 & 【講義20分】GitHub Pagesの仕組み、研究成果のWeb公開の意義。【演習70分】自分の研究紹介Webページの作成（GitHub Pages）。 & GitHub PagesによるWebページを公開し、URLを提出 \\
\hline
11 & 第9章前半 & ガクチカ・ES作成（基礎） & 【講義30分】就活スケジュール、企業の評価ポイント、STAR法の解説。【演習60分】AIを使った経験の棚卸しとSTAR法での構造化。 & STAR法で構造化した経験メモ（2テーマ分） \\
\hline
12 & 第9章後半 & ガクチカ・ES作成（実践） & 【講義20分】文章化のテクニック、面接との連動。【演習70分】ガクチカの文章化→AI添削→修正のサイクルを2回以上実施。ペアで相互レビュー。 & ガクチカ完成版（400字以内）、Before/After比較付き \\
\hline
13 & 第10章前半 & 研究の新規性分析 & 【講義30分】新規性の3つの観点、言語化の方法。【演習60分】自分の研究について新規性分析ワークフローを実施。 & 新規性分析レポート（テンプレートに沿って作成） \\
\hline
14 & 第10章後半 & 研究ポートフォリオの統合 & 【講義20分】ポートフォリオの構成、GitHubリポジトリの整理方法。【演習70分】これまでの成果物をGitHubリポジトリとして統合・整理。 & 研究ポートフォリオ（GitHubリポジトリ）の完成版を提出 \\
\hline
15 & 総合 & 最終発表と振り返り & 【発表60分】研究ポートフォリオのプレゼンテーション（5分×12名程度）。【振り返り30分】AI活用の全体振り返り、今後の活用方針の策定。 & AI活用方針書（A4用紙1枚）を最終課題として提出 \\
\end{longtable}

\section{授業運営のヒント}

\subsection*{事前準備}

\begin{itemize}
\item 受講者全員がChatGPTまたはClaudeのアカウントを持っていることを確認する（第1回の前に案内）
\item GitHubアカウントの作成も第1回の前に案内する
\item GitHub Copilotの学生無料枠の申請は手続きに時間がかかるため、早めに案内する
\end{itemize}

\subsection*{演習の進め方}

\begin{itemize}
\item 演習時間は、受講者が実際にAIと対話しながら作業する時間である。教員は巡回してサポートする
\item 可能であれば、受講者同士のペアワーク・グループワークを積極的に取り入れる。他の人のプロンプトや使い方を見ることが大きな学びとなる
\item 「うまくいかなかった例」も学びの素材として共有する
\end{itemize}

\subsection*{評価方法の案}

\begin{table}[H]
\centering
\small
\begin{tabular}{llp{6cm}}
\toprule
評価項目 & 配分 & 評価基準 \\
\midrule
各回の課題 & 40\% & 期限内の提出、指示への対応度 \\
中間発表 & 10\% & 研究内容の伝達力、スライドの質 \\
研究ポートフォリオ & 30\% & 成果物の質と統合度、GitHubの整理状況 \\
最終発表 & 10\% & プレゼンテーション力、質疑応答 \\
AI活用方針書 & 10\% & 振り返りの深さ、今後への具体的な方針 \\
\bottomrule
\end{tabular}
\end{table}

\subsection*{ポートフォリオの評価ルーブリック}

\begin{table}[H]
\centering
\footnotesize
\begin{tabular}{lp{1.8cm}p{1.8cm}p{1.8cm}p{1.8cm}}
\toprule
観点 & A（優秀） & B（良好） & C（合格） & D（要改善） \\
\midrule
成果物の網羅性 & 全章の成果物が揃っている & 大半の成果物が揃っている & 必須の成果物は揃っている & 成果物が不足している \\
\hline
成果物の質 & 高品質で研究内容が正確に反映 & 一定の質を持つ & 最低限の要件を満たしている & 質に問題がある \\
\hline
AI活用の適切さ & 効果的に活用しつつ批判的検証ができている & 概ね適切に検証している & AIを使っているが検証が不十分 & AIの出力を無批判に使用 \\
\hline
GitHubの整理 & README充実、構造明確、再現性確保 & 構造が概ね整理されている & 最低限の構造がある & 整理されていない \\
\hline
オリジナリティ & 独自の工夫や視点が随所に見られる & 一部に独自の工夫がある & テンプレート通り & 独自性がない \\
\bottomrule
\end{tabular}
\end{table}

\section{授業の変形パターン}

本書は15回の授業を想定しているが、以下のような変形も可能である。

\subsection*{8回集中講座の場合}

\begin{table}[H]
\centering
\begin{tabular}{cl}
\toprule
回 & 内容 \\
\midrule
1 & 第1章+第2章: AI基礎とプロンプトエンジニアリング \\
2 & 第3章+第4章: 文献調査と英語論文 \\
3 & 第5章: プレゼンテーション \\
4 & 第6章: プログラミング \\
5 & 第7章: データ分析 \\
6 & 第8章+第9章: Web公開と就活文書 \\
7 & 第10章: 新規性分析とポートフォリオ \\
8 & 最終発表 \\
\bottomrule
\end{tabular}
\end{table}

\subsection*{自主ゼミ（週1回、2時間）の場合}

各回の演習部分を充実させ、参加者同士のディスカッションを中心に進める。講義部分は事前に本書の該当章を読んできてもらい、演習とディスカッションに時間を割く反転授業形式が効果的である。
